\scanpage{9}
% slanted
\margin{L'énoncer \unsure{indifféremt|indifféremment} avec $\xi$ ou $L$}
\deleted{\unsure{$\underline L$ est|Lemme}} On va prouver ici que si $\xi\in\mathrm{Pic}(A/S)$, alors $\varphi_\xi=0$ \unsure{sss|ssi} on a
\[
\pi^*(\xi)-\mathrm{pr}_1^*(\xi)-\mathrm{pr}_2^*(\xi)=0
\]
En effet, considérons le diagramme
\diagram{\begin{tikzcd}[ampersand replacement=\&]
 A \& A\times A \arrow[l,"\mathrm{pr}_2"'] \arrow[r,"\pi"] \arrow[d,"\mathrm{pr}_1"] \& A\\
 \& A \&
\end{tikzcd}}
Dire que $\varphi_\xi=0$ signifie aussi que, après le chang$^{\mathrm{mt}}$ de base \deleted{\unsure{induit}} \unsure{sur|de|si} $S_1=A\unsure{\text{---}|\longrightarrow} S$, posant $A_1=A\times_S S_1$, et considérant la section universelle $t_1$ de $A_1$ sur $S_1$, on aura $t_1^*(\xi_1)=\xi_1$ (où $\xi_1=\xi\times_S S_1$).
% square at left of the lines beginning “posant”
\diagram{\begin{tikzcd}[ampersand replacement=\&]
A \arrow[d] \& A^{\unsure{\prime|\,}}_1=A\times_{\unsure{1|S}} A \arrow[l,"\mathrm{pr}_2"'] \arrow[d,"\mathrm{pr}_1"] \\
S \& A=S_1 \arrow[l,"\unsure{q|g}"]
\end{tikzcd}}
Or \deleted{$\unsure{t_1^*(\xi_1)}\,\unsure{t_1^*}(\xi_1)\unsure{=}\,t_1^*(\xi_1)={}$} Si $\xi$ est défini par $\underline L$ sur $A/S$, \inserted{trivial sur la section unité.} alors $\xi_1$ par $\mathrm{pr}_2^*(L)$, alors $t_1^*(\xi_1)$ est défini par
\[
\begin{aligned}
t_1^*(\mathrm{pr}_2^*(L))&=(\mathrm{pr}_2\,\mathrm{transl}(t_1))^*(\underline L)\\
&=\pi^*(\underline L).
\end{aligned}
\]
On doit donc \deleted{\unsure{prouver}} écrire que $\pi^*(\underline L)$ et $\mathrm{pr}_2^*(\underline L)$ définissent \unsure{la|le} $=$ \unsure{classe|élt} de $\mathrm{Pic}(A\times_S A/A)$, ou ce qui revient au même que
% slanted
\margin{l'écrire \unsure{tensoriellement}}
$\underbrace{\pi^*(\underline L)-\mathrm{pr}_1^*(\underline L)-\mathrm{pr}_2^*(\underline L)}$ définit \unsure{une|un} \unsure{classe|élt} \uline{\unsure{nulle|nul}} de $\mathrm{Pic}(A\times_S A/A)$. Comme ce faisceau induit un faisceau trivial sur la section unité de $A\times_S A$ sur $A$ :
\[
i_1^*(\text{---})=(\underbrace{\pi i_1}_{\mathrm{id}})^*(\underline L)-(\underbrace{\mathrm{pr}_1 i_1}_{\mathrm{id}})^*(\underline L)-(\underbrace{\mathrm{pr}_2 i_1}_{\varepsilon_{A/S}})^*(\underline L)=0
\]
% archive sheet number: 8

